
\chapter{TỔNG KẾT VÀ HƯỚNG PHÁT TRIỂN}

% ============================================
% 6.1. TỔNG KẾT QUÁ TRÌNH THỰC HIỆN
% ============================================

\section{Tổng kết quá trình thực hiện}

Trong khuôn khổ học phần Nhập môn Học máy, nhóm đã
thực hiện đề tài ``Phân loại cảm xúc từ phản hồi của
sinh viên'' với mục tiêu xây dựng và đánh giá một
quy trình học máy có giám sát nhằm phân loại
cảm xúc từ dữ liệu văn bản.

Quá trình thực hiện được triển khai theo các giai
đoạn chính: khám phá dữ liệu, tiền xử lý văn bản,
biểu diễn đặc trưng, huấn luyện mô hình, đánh giá
kết quả và phân tích những trường hợp dự đoán sai.

Nhóm sử dụng bộ dữ liệu Contextual Sentiment
Feedback Dataset được giảng viên chỉ định,
gồm 1.223 phản hồi và 11 thuộc tính. Bài toán
được xác định dưới dạng phân loại văn bản
thành ba lớp cảm xúc: positive, neutral
và negative.

Ở giai đoạn khám phá dữ liệu, nhóm tiến hành
kiểm tra cấu trúc, giá trị thiếu, bản ghi
trùng lặp, phân phối nhãn và độ dài văn bản.
Kết quả cho thấy bộ dữ liệu không có giá trị
thiếu hoặc phản hồi trùng lặp theo các phép
kiểm tra đã thực hiện. Tuy nhiên, một số
phản hồi có nội dung ngữ nghĩa chưa rõ ràng,
cần được chú ý khi phân tích kết quả.

Ở giai đoạn tiền xử lý, nhóm thực hiện chuẩn hóa
chữ thường, xử lý URL, dấu câu, ký tự đặc biệt
và khoảng trắng thừa. Các từ phủ định và những
thành phần có khả năng chứa thông tin cảm xúc
được giữ lại nhằm hạn chế việc loại bỏ thông
tin một cách máy móc.

Sau khi làm sạch, dữ liệu được chia thành
978 mẫu huấn luyện và 245 mẫu kiểm thử,
tương ứng tỷ lệ 80/20, với
\texttt{random\_state=42} và phương pháp
phân tầng theo nhãn.

Nhóm xây dựng bốn cấu hình biểu diễn văn bản
dựa trên Bag-of-Words và TF-IDF, kết hợp với
unigram hoặc unigram--bigram.

Trên cơ sở đó, hai thuật toán Multinomial
Naive Bayes và K-Nearest Neighbors (K-NN)
được sử dụng để xây dựng các mô hình phân
loại và thực hiện thử nghiệm so sánh.

Quá trình lựa chọn siêu tham số được thực
hiện bằng Stratified 3-Fold Cross-Validation
trên tập huấn luyện với Macro F1-score
làm tiêu chí đánh giá.

Nhóm cũng xây dựng phương án baseline
để làm mốc so sánh và sử dụng các chỉ số
Accuracy, Precision, Recall, F1-score
cùng Confusion Matrix để đánh giá
kết quả cuối cùng.

Toàn bộ quy trình được tổ chức thành bốn
Jupyter Notebook, tương ứng với các giai
đoạn khám phá dữ liệu, tiền xử lý, huấn luyện
và đánh giá mô hình.


% ============================================
% 6.2. KẾT QUẢ ĐẠT ĐƯỢC
% ============================================

\section{Kết quả đạt được}

\subsection{Kết quả về xây dựng quy trình học máy}

Thông qua quá trình thực hiện, nhóm đã xây dựng
được một quy trình phân loại văn bản bao gồm
các thành phần cần thiết của bài toán học máy
có giám sát.

Các kết quả chính bao gồm:

\begin{itemize}

    \item Hoàn thành việc khám phá và phân tích
    bộ dữ liệu phản hồi sinh viên, xác định
    cấu trúc dữ liệu, nhãn mục tiêu, phân phối
    lớp và những vấn đề cần lưu ý về chất
    lượng văn bản.

    \item Xây dựng quy trình tiền xử lý văn bản
    và kiểm tra dữ liệu trước khi huấn luyện.

    \item Triển khai hai phương pháp biểu diễn
    văn bản Bag-of-Words và TF-IDF với bốn
    cấu hình unigram và unigram--bigram.

    \item Huấn luyện và đánh giá hai thuật toán
    Multinomial Naive Bayes và K-NN trên cùng
    bộ dữ liệu và quy trình đánh giá.

    \item Thực hiện 24 lượt thử siêu tham số
    bằng Stratified 3-Fold Cross-Validation
    trên tập huấn luyện.

    \item Xây dựng baseline và so sánh các
    phương án thử nghiệm dựa trên nhiều
    chỉ số đánh giá.

    \item Tạo các bảng kết quả, biểu đồ
    so sánh, Confusion Matrix và danh sách
    phản hồi dự đoán sai để phục vụ
    phân tích thực nghiệm.

\end{itemize}

Việc tách dữ liệu Train/Test trước khi học
bộ đặc trưng, kết hợp với việc sử dụng
Pipeline trong Cross-Validation, giúp hạn
chế nguy cơ rò rỉ dữ liệu trong quá trình
lựa chọn mô hình.

Bên cạnh đó, nhóm đã lưu lại các tệp dữ liệu
trung gian, cấu hình thử nghiệm và kết quả
dự đoán nhằm hỗ trợ kiểm tra, đối chiếu
và thực hiện lại quy trình.


\subsection{Kết quả thực nghiệm}

Kết quả kiểm định chéo cho thấy mô hình
TF-IDF Unigram kết hợp K-NN với $k=5$
đạt Macro F1-score trung bình cao nhất
trong các cấu hình được thử nghiệm,
với giá trị 0,3647.

Tuy nhiên, khi đánh giá trên tập kiểm thử
gồm 245 phản hồi, mô hình được lựa chọn
chỉ đạt Accuracy 0,3388 và Macro F1-score
0,3369.

Trong khi đó, mô hình baseline sử dụng
Bag-of-Words Unigram kết hợp Multinomial
Naive Bayes với $\alpha=1,0$ đạt
Accuracy 0,3878 và Macro F1-score 0,3748.

Như vậy, Macro F1-score trên tập kiểm thử
của mô hình được lựa chọn bằng Cross-Validation
thấp hơn baseline 0,0379.

Kết quả này cho thấy phương án có điểm
kiểm định chéo cao nhất chưa đạt hiệu quả
tốt hơn baseline trên tập kiểm thử hiện tại.

Thông qua Confusion Matrix, nhóm xác định
mô hình được lựa chọn dự đoán đúng 83
trong tổng số 245 phản hồi kiểm thử
và dự đoán sai 162 phản hồi.

Cặp nhãn có số trường hợp nhầm lẫn
nhiều nhất là positive được dự đoán
thành negative, với 41 trường hợp.

Các kết quả trên cho thấy mô hình vẫn
gặp khó khăn trong việc phân biệt
ba nhóm cảm xúc.

Mặc dù hiệu quả phân loại chưa cao,
quá trình thực nghiệm đã cung cấp
những kết quả cụ thể để nhóm so sánh
các phương pháp, phân tích hạn chế
và đề xuất hướng cải tiến.

Các bảng kết quả và ma trận nhầm lẫn
chi tiết đã được trình bày tại
Chương 4 và thảo luận trong Chương 5.


% ============================================
% 6.3. ĐÁNH GIÁ MỨC ĐỘ HOÀN THÀNH
% ============================================

\section{Đánh giá mức độ hoàn thành mục tiêu}

Đối chiếu với các mục tiêu đã xác định
trong Chương 1, nhóm nhận thấy những
nội dung chính về xây dựng và thực nghiệm
quy trình học máy đã được triển khai.

Thứ nhất, nhóm đã hoàn thành việc khám phá,
kiểm tra và phân tích các đặc điểm cơ bản
của bộ dữ liệu được giảng viên cung cấp.

Thứ hai, nhóm đã xây dựng được quy trình
tiền xử lý văn bản và chuyển đổi văn bản
thành biểu diễn số thông qua Bag-of-Words
và TF-IDF.

Thứ ba, nhóm đã triển khai hai thuật toán
phân loại đã học là Naive Bayes và K-NN,
đồng thời thực hiện thử nghiệm thay đổi
siêu tham số và so sánh nhiều cấu hình.

Thứ tư, nhóm đã đánh giá mô hình bằng
các chỉ số Accuracy, Precision, Recall,
F1-score và Confusion Matrix, thay vì
chỉ sử dụng một giá trị độ chính xác.

Thứ năm, nhóm đã xem xét các trường hợp
dự đoán sai và thảo luận những hạn chế
có thể liên quan đến chất lượng dữ liệu,
phương pháp biểu diễn văn bản và
khả năng khái quát hóa của mô hình.

Như vậy, các mục tiêu về triển khai
quy trình, so sánh thuật toán và
phân tích kết quả đã được thực hiện
trong phạm vi bài tập lớn.

Tuy nhiên, cần phân biệt giữa việc hoàn
thành quy trình thực nghiệm và việc đạt
được một mô hình có chất lượng dự đoán cao.

Kết quả hiện tại chưa chứng minh được
mô hình có độ tin cậy đủ tốt để sử dụng
trực tiếp trong các hoạt động phân tích
phản hồi sinh viên ngoài thực tế.

Do đó, nhóm đánh giá rằng mục tiêu
xây dựng và nghiên cứu quy trình
học máy cơ bản đã được thực hiện,
trong khi mục tiêu nâng cao hiệu quả
và khả năng ứng dụng thực tế vẫn cần
được tiếp tục cải thiện.


% ============================================
% 6.4. HẠN CHẾ VÀ HƯỚNG PHÁT TRIỂN
% ============================================

\section{Hạn chế và hướng phát triển}

\subsection{Những hạn chế còn tồn tại}

Mặc dù đã triển khai đầy đủ các giai đoạn
chính của quy trình học máy, kết quả thực
nghiệm vẫn tồn tại một số hạn chế.

Trước hết, bộ dữ liệu có quy mô còn giới hạn
và một số phản hồi có nội dung thiếu rõ
nghĩa. Điều này đặt ra yêu cầu cần tiếp tục
xem xét tính phù hợp giữa nội dung văn bản
và nhãn cảm xúc được cung cấp.

Tiếp theo, hai phương pháp Bag-of-Words
và TF-IDF chủ yếu dựa trên sự xuất hiện
và trọng số của từ hoặc cụm từ.

Vì vậy, các phương pháp này chưa biểu diễn
đầy đủ ngữ cảnh và quan hệ ngữ nghĩa phức
tạp giữa các thành phần trong câu.

Bên cạnh đó, phạm vi thử nghiệm siêu tham
số còn tương đối nhỏ. Nhóm mới khảo sát
ba giá trị $\alpha$ đối với Naive Bayes
và ba giá trị $k$ đối với K-NN.

Quá trình đánh giá cũng mới sử dụng
một lần chia Train/Test cố định.
Do đó, kết quả chưa phản ánh đầy đủ
mức độ biến thiên có thể xảy ra
khi thay đổi cách phân chia dữ liệu.

Cuối cùng, hiệu quả phân loại trên tập
kiểm thử còn thấp, đặc biệt đối với
khả năng phân biệt các phản hồi
thuộc ba lớp cảm xúc.

Những hạn chế trên cho thấy cần
tiếp tục kiểm chứng và cải thiện
giải pháp trước khi xem xét
khả năng ứng dụng thực tế.


\subsection{Hướng cải thiện chất lượng dữ liệu}

Một trong những hướng phát triển cần
được ưu tiên là kiểm tra và nâng cao
chất lượng dữ liệu đầu vào.

Nhóm có thể xây dựng bộ tiêu chí
đánh giá nhãn cảm xúc và thực hiện
kiểm tra thủ công một số lượng mẫu
đại diện cho từng lớp.

Các phản hồi có nội dung khó hiểu,
thiếu ngữ cảnh hoặc có dấu hiệu
không thống nhất giữa văn bản
và nhãn cần được xem xét kỹ hơn.

Trong trường hợp có điều kiện,
nhóm có thể nghiên cứu sử dụng
thêm dữ liệu phản hồi được thu
thập và gán nhãn theo một quy
trình thống nhất.

Việc mở rộng hoặc thay đổi bộ
dữ liệu trong phạm vi bài tập lớn
cần tuân thủ yêu cầu và được sự
đồng ý của giảng viên.


\subsection{Hướng cải thiện tiền xử lý và biểu diễn đặc trưng}

Trong các nghiên cứu tiếp theo,
nhóm có thể thực hiện thêm các
thí nghiệm nhằm xem xét tác động
của những phương pháp tiền xử lý
khác nhau.

Chẳng hạn, nhóm có thể so sánh
việc giữ và loại bỏ một số dấu câu,
điều chỉnh các ngưỡng
\texttt{min\_df}, \texttt{max\_df}
hoặc số lượng đặc trưng tối đa.

Bên cạnh đó, nhóm có thể tiếp tục
khảo sát những cấu hình n-gram
khác nhau để đánh giá khả năng
biểu diễn các cụm từ mang
thông tin cảm xúc.

Nếu có thêm thời gian và tài nguyên
tính toán, nhóm có thể tìm hiểu
các phương pháp biểu diễn văn bản
nâng cao, chẳng hạn như
Word Embeddings hoặc Sentence Embeddings.

Tuy nhiên, những phương pháp mở rộng
cần được lựa chọn có lý do, thực hiện
thí nghiệm so sánh và đánh giá
trên dữ liệu phù hợp trước khi
đưa ra kết luận về hiệu quả.


\subsection{Hướng cải thiện mô hình và đánh giá}

Đối với giai đoạn huấn luyện,
nhóm có thể mở rộng phạm vi
tìm kiếm các giá trị siêu tham số
của Naive Bayes và K-NN.

Ngoài ra, việc thực hiện kiểm định
chéo lặp lại hoặc sử dụng nhiều
lần chia dữ liệu phù hợp có thể
giúp nhóm xem xét mức độ ổn định
của kết quả.

Các thí nghiệm bổ sung cần tiếp tục
tuân thủ nguyên tắc lựa chọn tham
số trên dữ liệu huấn luyện hoặc
validation, không sử dụng tập
test để điều chỉnh mô hình.

Trong trường hợp có thêm dữ liệu,
nhóm có thể xây dựng một tập
kiểm thử độc lập mới để đánh giá
khả năng khái quát hóa của giải pháp.

Bên cạnh các chỉ số tổng hợp,
cần tiếp tục phân tích hiệu quả
theo từng lớp cảm xúc và kiểm tra
các trường hợp nhầm lẫn quan trọng.


\subsection{Khả năng phát triển ứng dụng}

Nếu mô hình được cải thiện và
đánh giá đạt chất lượng phù hợp,
nhóm có thể xem xét xây dựng
một ứng dụng thử nghiệm cho phép
người dùng nhập nội dung phản hồi
và nhận kết quả phân loại cảm xúc.

Ứng dụng có thể hỗ trợ thống kê
số lượng phản hồi theo các nhóm
tích cực, trung tính và tiêu cực.

Tuy nhiên, kết quả phân loại cần
được sử dụng như thông tin tham
khảo, đặc biệt đối với các phản
hồi có ngữ cảnh phức tạp hoặc
không thể hiện rõ cảm xúc.

Trong trường hợp triển khai thực tế,
cần bổ sung các biện pháp bảo vệ
dữ liệu cá nhân, kiểm soát quyền
truy cập và cơ chế kiểm tra thủ
công đối với những dự đoán
không chắc chắn.

Đây là hướng phát triển trong
tương lai, không phải chức năng
đã được triển khai trong phạm
vi thực nghiệm hiện tại.


% ============================================
% 6.5. KẾT LUẬN CHUNG
% ============================================

\section{Kết luận chung}

Thông qua đề tài ``Phân loại cảm xúc từ phản hồi
của sinh viên'', nhóm đã vận dụng những kiến thức
cơ bản của học phần Nhập môn Học máy để xây dựng
và thực nghiệm một quy trình phân loại văn bản
có giám sát.

Đề tài đã thực hiện các bước từ khám phá dữ liệu,
tiền xử lý, biểu diễn đặc trưng, huấn luyện
mô hình đến đánh giá và phân tích kết quả.

Việc sử dụng hai thuật toán Multinomial Naive Bayes
và K-NN, kết hợp với Bag-of-Words và TF-IDF,
giúp nhóm có cơ sở so sánh ảnh hưởng của phương
pháp biểu diễn và các tham số đối với hiệu quả
phân loại.

Kết quả thực nghiệm cho thấy mô hình được lựa
chọn bằng Cross-Validation chưa đạt hiệu quả
tốt hơn baseline trên tập kiểm thử hiện tại.

Đồng thời, quá trình phân tích lỗi đã chỉ ra
những khó khăn trong việc phân biệt các lớp
cảm xúc, đặc biệt với những phản hồi có
nội dung thiếu ngữ cảnh hoặc chưa rõ nghĩa.

Từ kết quả này, nhóm nhận thấy hiệu quả của
một hệ thống học máy không chỉ phụ thuộc vào
thuật toán mà còn liên quan chặt chẽ đến chất
lượng dữ liệu, cách biểu diễn đặc trưng và
phương pháp đánh giá.

Mặc dù kết quả phân loại còn hạn chế,
đề tài đã giúp nhóm thực hành các bước
quan trọng trong việc xây dựng một
quy trình học máy hoàn chỉnh, đồng thời
rèn luyện khả năng phân tích dữ liệu,
thiết kế thực nghiệm và diễn giải kết quả.

Những kinh nghiệm thu được là cơ sở để
nhóm tiếp tục nghiên cứu, cải thiện
mô hình và phát triển các ứng dụng
phân tích văn bản trong tương lai.
